\section{研究背景与因子简介}
单因子测试是多因子研究的起点，其作用是检验一个信号能否在同一期股票之间区分后续收益，并判断这种关系的稳定性和适用范围。本报告按照课程附录，以回归、Rank IC与等权分层三类证据评价BP，再考察年度、行业和市值域的差异。研究问题是：控制行业与规模后，较高BP是否仍对应较高下期收益，以及这一关系在哪些范围内失效。

BP属于估值类因子。较高BP意味着账面权益相对市场定价较高，可能反映较低估值，也可能包含较高经营风险。本文先给出操作定义，再检验正向假说；预期方向在计算绩效之前确定。

\begin{table}[H]\centering
\caption*{\textbf{表1　因子定义与默认方向}}
\small\renewcommand{\arraystretch}{1.2}
\begin{tabularx}{\linewidth}{@{}p{26mm}X@{}}\toprule
项目 & 定义 \\ \midrule
因子 & 账面市值比（Book-to-Price，BP） \\
经济概念 & 账面权益与市场价值之比 \\
实际计算 & $\mathrm{BP}_{i,t}=1/\mathrm{PB}_{i,t}$，仅使用正且有限的历史PB \\
源字段 & 课程离线估值表的\texttt{pb}；供应商当前定义为总市值除以“净资产减其他权益工具” \\
单位与频率 & 无量纲比值；交易月末收盘后取值 \\
预期方向 & 正向。Q1为低中性化残差层，Q5为高中性化残差层；十层时Q10最高 \\
无效值 & 缺失、零、负和无穷PB均不填补；这是定义域限制，不属于去极值 \\
控制与分域 & 控制规模使用对数总市值；WLS权重和市值分域使用流通市值，源市值单位为万元 \\
\bottomrule\end{tabularx}
\end{table}

实际检验依赖供应商历史日度PB，不能直接等同于已自行重构的、当时可获得且未经修订的账面权益序列。已有字段定义和有限财务样本支持这一操作口径，但严格的逐版本时点可得性（point-in-time，PIT）尚不能证明。

高BP后续取得较高收益，可能与错误定价修正或风险补偿有关；资产减值、盈利能力下降和持续经营风险也可能构成“价值陷阱”。这些解释并不互斥，本文的观察性检验不能识别其因果机制。非线性和适用域分析用于限定结论，而非根据结果改选因子或反转方向。
